\section{泛定方程与定解问题}
\idx{泛定方程}{fandingfangcheng}
\idx{定解条件}{dingjietiaojian}
\idx{定解问题}{dingjiewenti}
\idx{初始条件}{chushitiaojian}
\idx{边界条件}{bianjietiaojian}
\idx{衔接条件}{xianjietiaojian}

未知量在区域内部必须满足的偏微分方程,称为\textbf{泛定方程}.
上一节导出的波动、热传导、拉普拉斯与泊松方程都是泛定方程,写的是这类问题的共性.
代入后能使方程变成恒等式的函数,称为该泛定方程的解.
这样的解一般不唯一:
细弦的小横振动都满足同一条振动方程,
轻轻拨一下和在中点敲一下都是它的解,
差别在开始时的位置与速度,以及两端的装置.

用来从这些解里选出符合这一次具体过程的条件,总称\textbf{定解条件}.
\textbf{初始条件}给出某一时刻的状态.
对时间的导数有多高,就要给多少项;方程里没有时间,就没有初始条件.
\textbf{边界条件}写在区域边界上.
第一类给定未知量本身,第二类给定法向导数,第三类给定二者的线性组合.
右端为零称为齐次,右端为已知非零函数则称为非齐次.
\textbf{衔接条件}用在区域内部介质突然改变的分界面上,不是边界条件.
通常要求未知量连续,以及两侧的力或通量互相平衡.

一条泛定方程,加上为这次过程配好的定解条件,构成一个\textbf{定解问题}.
所求的函数在区域内部满足方程,在初始时刻和边界上满足相应条件;
介质内部有接点时,还要满足衔接条件.

若解存在、唯一,并且定解数据作小改动时解也只作小改动,则称该定解问题\textbf{适定}.
条件不够则解不唯一,条件互相矛盾则解可能不存在;这两种情形都不讨论.
以后只讨论适定的定解问题.
泛定方程和定解条件都是线性的时候,适用\textbf{叠加原理}:
若 $u_1$、$u_2$ 分别满足数据为 $f_1,\varphi_1,\ldots$ 与 $f_2,\varphi_2,\ldots$ 的同一类定解问题,
则 $c_1 u_1+c_2 u_2$ 满足数据为相应线性组合的问题.

\subsection{弦振动问题的定解条件}

上一节已导出弦振动的泛定方程.它只给出位移 $u(x,t)$ 所满足的一般规律,
还不能单独确定具体弦的运动;还须补上定解条件.
下面按梁昆淼书的提法,把三类边界条件与初始条件写到弦振动问题上,并组成完整的定解问题.

若弦的两端固定在 $x=0$ 与 $x=l$,则
\begin{equation}
    u(0, t)=0, \quad u(l, t)=0,
    \label{eq:boundary_conditions}
\end{equation}
这是\textbf{第一类边界条件}（Dirichlet）.
设初始时刻的位移与速度为
\begin{equation}
    u(x, 0)=\varphi(x), \quad \frac{\partial u(x, 0)}{\partial t}=\psi(x) \quad(0 \leqslant x \leqslant l),
    \label{eq:initial_conditions}
\end{equation}
这是\textbf{初始条件}.
把泛定方程与上述定解条件合在一起,得到定解问题
$$
\left\{\begin{array}{l}
	\dfrac{\partial^2 u}{\partial t^2}-a^2 \dfrac{\partial^2 u}{\partial x^2}=f(x, t), \\
	t=0: u=\varphi(x), \dfrac{\partial u}{\partial t}=\psi(x), \\
	x=0: \quad u=0, \\
	x=l: \quad u=0 .
\end{array}\right.
$$
即要求 $u(x,t)$ 在 $0<x<l$, $t>0$ 内满足泛定方程,
在 $t=0$ 上满足\eqref{eq:initial_conditions},
在 $x=0$ 与 $x=l$ 上满足\eqref{eq:boundary_conditions}.

弦振动的边界条件还可以有以下两种:

(a) 弦的一端 (例如 $x=0$ ) 处于自由状态, 即可以在垂直于 $x$ 轴的直线上自由滑动, 未受到垂直方向外力.
在边界右端的张力的垂直方向分量是 $T \frac{\partial u}{\partial x}$, 得出此时应成立
$$
\left.\frac{\partial u}{\partial x}\right|_{x=0}=0 .
$$
也可以考虑更普遍的边界条件
$$
\left.\frac{\partial u}{\partial x}\right|_{x=0}=\mu(t) ,
$$
其中 $\mu(t)$ 是 $t$ 的已知函数.这种边界条件称为\textbf{第二类边界条件} (又称诺伊曼 (Neumann) 边界条件).

(b)在应用上还会遇到另一种情形.将弦的一端固定在弹性支承上,也就是说此时支承的伸缩符合胡克定律. 
如果支承原来的位置为 $u=0$, 则 $u$ 在端点的值表示支承在该点的伸长.例如在 $x=l$ 的一端, 
弦对支承拉力的垂直方向分量为 $-T \frac{\partial u}{\partial x}$, 由胡克定律知
$$
-\left.T \frac{\partial u}{\partial x}\right|_{x=l}=\left.k u\right|_{x=l},
$$

其中 $k$ 为弹性系数. 因此在弹性支承的情形,边界条件归结为
$$
\left.\left(\frac{\partial u}{\partial x}+\sigma u\right)\right|_{x=l}=0
$$
其中 $\sigma=\frac{k}{T}$ 是已知正数.在数学中也可以考虑更普遍的边界条件
$$
\left.\left(\frac{\partial u}{\partial x}+\sigma u\right)\right|_{x=l}=v(t),
$$
其中 $v(t)$ 是 $t$ 的已知函数.这种边界条件称为\textbf{第三类边界条件}.

\eqref{eq:boundary_conditions} 是齐次的第一类边界条件;
若改为 $u|_{x=0}=\mu_1(t)$, $u|_{x=l}=\mu_2(t)$, 则为非齐次第一类边界条件.
初始条件\eqref{eq:initial_conditions} 在 $\varphi,\psi$ 不全为零时是非齐次的;
$\varphi\equiv\psi\equiv 0$ 时是齐次初始条件.

\subsection{热传导与稳定场的定解条件}

热传导方程对时间是一阶,初始条件只给初始温度 $u|_{t=0}=\varphi$.
边界仍是三类:给定表面温度;给定法向导数 $\partial u/\partial n$(绝热时为零);
或给定温度与法向导数的线性组合(表面与外界换热).

稳定场方程不含时间,没有初始条件,只由边界条件确定解.
导体表面电势给定,就是第一类边界条件.

介质在区域内部突然改变时,分界面上还要写衔接条件.
设接点为 $x=\xi$,两侧无集中质量、无面热源.
两段弦相接时位移连续且横向力平衡,
$u|_{\xi^-}=u|_{\xi^+}$,
$T_-\partial_x u|_{\xi^-}=T_+\partial_x u|_{\xi^+}$.
两段导热杆相接时温度连续且热流连续,
$u|_{\xi^-}=u|_{\xi^+}$,
$k_-\partial_x u|_{\xi^-}=k_+\partial_x u|_{\xi^+}$.
接点上若另有集中质量或面热源,第二个等式改为两侧之差等于该集中量.

